% -----------------------------------------------------------------------------
% Numerical Methods lecture notes preamble — circuit notebook style
% Structure: thmtools + mdframed
% Compile with a class that has chapters, e.g. \documentclass{report} or book.
% -----------------------------------------------------------------------------

% Basic commands
\usepackage[margin=1.45in]{geometry}
\usepackage[spanish, es-noquoting]{babel}
\usepackage[T1]{fontenc}
\usepackage[utf8]{inputenc}
\usepackage{lmodern}
\usepackage{microtype}
\usepackage{amsmath, amsfonts, mathtools, amsthm, amssymb}
\usepackage[usenames,dvipsnames,svgnames]{xcolor}
\usepackage{float}
\usepackage{enumitem}
\setlist{itemsep=0.15em, topsep=0.35em}

% TikZ and diagrams
\usepackage{tikz}
\usepackage{tikz-cd}
\usetikzlibrary{intersections, angles, quotes, calc, positioning}
\usetikzlibrary{arrows.meta, decorations.pathmorphing, shadows.blur}
\usepackage{pgfplots}
\pgfplotsset{compat=1.18}

% Hyperlinks
\usepackage[colorlinks=true, linkcolor=black, citecolor=black, urlcolor=black]{hyperref}

% Theorem machinery
\usepackage{thmtools}
\usepackage[framemethod=TikZ]{mdframed}
\mdfsetup{skipabove=1em, skipbelow=0.2em}

% Useful math commands
\newcommand{\R}{\mathbb{R}}
\newcommand{\C}{\mathbb{C}}
\newcommand{\F}{\mathbb{F}}
\newcommand{\N}{\mathbb{N}}
\newcommand{\Q}{\mathbb{Q}}
\newcommand{\Z}{\mathbb{Z}}
\newcommand{\K}{\mathbb{K}}
\newcommand{\eps}{\varepsilon}
\let\epsilon\varepsilon
\newcommand{\ds}{\displaystyle}
\newcommand{\ol}{\overline}
\newcommand{\ul}{\underline}
\newcommand{\xto}[1]{\xrightarrow{#1}}
\newcommand{\xfrom}[1]{\xleftarrow{#1}}
\DeclareMathOperator{\Ker}{Ker}
\DeclareMathOperator{\Imagen}{Im}
\DeclareMathOperator{\Hom}{Hom}
\DeclareMathOperator{\Aut}{Aut}
\DeclareMathOperator{\End}{End}
\DeclareMathOperator{\GL}{GL}
\DeclareMathOperator{\SL}{SL}
\DeclareMathOperator{\SO}{SO}
\DeclareMathOperator{\grad}{grad}
\DeclareMathOperator{\traz}{traz}
\DeclareMathOperator{\Int}{Int}
\DeclareMathOperator{\cl}{cl}
\DeclareMathOperator{\id}{id}
\DeclareMathOperator{\dist}{dist}
\DeclareMathOperator{\rank}{rank}
\DeclareMathOperator{\Span}{span}
\DeclareMathOperator{\diam}{diam}
\DeclareMathOperator{\supp}{supp}
\DeclareMathOperator{\Spec}{Spec}
\DeclareMathOperator{\Var}{Var}
\DeclareMathOperator{\Cov}{Cov}
\DeclareMathOperator{\ord}{ord}
\DeclareMathOperator{\dom}{dom}
\DeclareMathOperator{\ran}{ran}
\DeclareMathOperator{\mcd}{mcd}
\DeclareMathOperator{\mcm}{mcm}

% Header/footer
\usepackage{fancyhdr}
\pagestyle{fancy}
\fancyhf{}
\fancyhead[L]{\sffamily\nouppercase{Victoria Eugenia Torroja}}
\fancyhead[R]{\sffamily\nouppercase\rightmark}
\fancyfoot[C]{\sffamily\leftmark}
\fancyfoot[R]{\sffamily\thepage}
\renewcommand{\headrulewidth}{0.4pt}
\renewcommand{\footrulewidth}{0pt}

% Generic proof name in Spanish
\renewcommand{\proofname}{Demostración}

% -----------------------------------------------------------------------------
% Numerical Methods — circuit notebook style
% -----------------------------------------------------------------------------

% Palette: technical orange/blue on clean paper
\definecolor{NumInk}{HTML}{1F2937}
\definecolor{NumBlue}{HTML}{1E5A88}
\definecolor{NumOrange}{HTML}{C45A12}
\definecolor{NumPaper}{HTML}{F7FAFC}
\definecolor{NumCream}{HTML}{FFF7EA}

\colorlet{ResultTitle}{NumBlue!85!black}
\colorlet{DefinitionTitle}{NumOrange!85!black}
\colorlet{ExampleTitle}{NumInk}
\colorlet{NotationTitle}{NumBlue!75!black}

% mdframed styles
\mdfdefinestyle{numresultframe}{%
    nobreak,
    roundcorner=0pt,
    linewidth=0.8pt,
    linecolor=NumBlue,
    backgroundcolor=NumPaper,
    leftmargin=0pt,
    rightmargin=0pt,
    innerleftmargin=9pt,
    innerrightmargin=9pt,
    innertopmargin=7pt,
    innerbottommargin=8pt,
    skipabove=0.85em,
    skipbelow=0.85em,
    shadow=true,
    shadowsize=4pt,
    shadowcolor=black!16,
    splittopskip=1.2em,
    splitbottomskip=0.8em,
    firstextra={\draw[NumBlue, line width=3pt] (P) -- (P-|O);},
    middleextra={\draw[NumBlue, line width=3pt] (P) -- (P-|O);},
    secondextra={\draw[NumBlue, line width=3pt] (P) -- (P-|O);},
    singleextra={\draw[NumBlue, line width=3pt] (P) -- (P-|O);}
}

\mdfdefinestyle{numdefinitionframe}{%
    nobreak,
    roundcorner=0pt,
    linewidth=0.8pt,
    linecolor=NumOrange,
    backgroundcolor=NumCream,
    leftmargin=0pt,
    rightmargin=0pt,
    innerleftmargin=9pt,
    innerrightmargin=9pt,
    innertopmargin=7pt,
    innerbottommargin=8pt,
    skipabove=0.85em,
    skipbelow=0.85em,
    shadow=true,
    shadowsize=4pt,
    shadowcolor=black!16,
    splittopskip=1.2em,
    splitbottomskip=0.8em,
    firstextra={\draw[NumOrange, line width=3pt] (o) -- (P-|O);},
    middleextra={\draw[NumOrange, line width=3pt] (P) -- (P-|O);},
    secondextra={\draw[NumOrange, line width=3pt] (P) -- (P-|O);},
    singleextra={\draw[NumOrange, line width=3pt] (P) -- (P-|O);}
}

\mdfdefinestyle{numexampleframe}{%
    roundcorner=0pt,
    linewidth=0.8pt,
    linecolor=NumInk!60,
    backgroundcolor=white,
    leftmargin=0pt,
    rightmargin=0pt,
    innerleftmargin=9pt,
    innerrightmargin=9pt,
    innertopmargin=7pt,
    innerbottommargin=8pt,
    skipabove=0.85em,
    skipbelow=0.85em,
    shadow=true,
    shadowsize=3pt,
    shadowcolor=black!9,
    splittopskip=1.2em,
    splitbottomskip=0.8em,
    firstextra={\draw[NumInk!60, line width=2.5pt] (O) -- (P-|O);},
    middleextra={\draw[NumInk!60, line width=2.5pt] (O) -- (P-|O);},
    secondextra={\draw[NumInk!60, line width=2.5pt] (O) -- (P-|O);},
    singleextra={\draw[NumInk!60, line width=2.5pt] (O) -- (P-|O);}
}

\mdfdefinestyle{numobsframe}{%
    roundcorner=0pt,
    linewidth=0.8pt,
    linecolor=Maroon!60,
    backgroundcolor=white,
    leftmargin=0pt,
    rightmargin=0pt,
    innerleftmargin=9pt,
    innerrightmargin=9pt,
    innertopmargin=7pt,
    innerbottommargin=8pt,
    skipabove=0.85em,
    skipbelow=0.85em,
    shadow=true,
    shadowsize=3pt,
    shadowcolor=black!9,
    splittopskip=1.2em,
    splitbottomskip=0.8em,
    firstextra={\draw[Maroon!60, line width=2.5pt] (O) -- (P-|O);},
    middleextra={\draw[Maroon!60, line width=2.5pt] (O) -- (P-|O);},
    secondextra={\draw[Maroon!60, line width=2.5pt] (O) -- (P-|O);},
    singleextra={\draw[Maroon!60, line width=2.5pt] (O) -- (P-|O);}
}

\mdfdefinestyle{numnotationframe}{%
    roundcorner=0pt,
    linewidth=0.8pt,
    linecolor=NumBlue!55,
    backgroundcolor=white,
    leftmargin=0pt,
    rightmargin=0pt,
    innerleftmargin=9pt,
    innerrightmargin=9pt,
    innertopmargin=7pt,
    innerbottommargin=8pt,
    skipabove=0.85em,
    skipbelow=0.85em,
    shadow=true,
    shadowsize=3pt,
    shadowcolor=black!9,
    splittopskip=1.2em,
    splitbottomskip=0.8em,
    firstextra={\draw[NumBlue!55, line width=2.5pt] (O) -- (P-|O);},
    middleextra={\draw[NumBlue!55, line width=2.5pt] (O) -- (P-|O);},
    secondextra={\draw[NumBlue!55, line width=2.5pt] (O) -- (P-|O);},
    singleextra={\draw[NumBlue!55, line width=2.5pt] (O) -- (P-|O);}
}

% Numerical methods helpers
\DeclareMathOperator{\cond}{cond}
\DeclareMathOperator{\fl}{fl}
\DeclareMathOperator{\diag}{diag}
\DeclareMathOperator{\tol}{tol}
\newcommand{\Oh}{\mathcal{O}}
\newcommand{\norm}[1]{\left\lVert #1\right\rVert}
\newcommand{\abs}[1]{\left\lvert #1\right\rvert}

% Theorem styles
\declaretheoremstyle[
    headfont=\bfseries\sffamily\color{DefinitionTitle},
    bodyfont=\normalfont,
    spaceabove=0pt,
    spacebelow=0pt,
    mdframed={style=numdefinitionframe}
]{numdefinitionbox}

\declaretheoremstyle[
    headfont=\bfseries\sffamily\color{ResultTitle},
    bodyfont=\normalfont,
    spaceabove=0pt,
    spacebelow=0pt,
    mdframed={style=numresultframe}
]{numresultbox}

\declaretheoremstyle[
    headfont=\bfseries\sffamily\color{ExampleTitle},
    bodyfont=\normalfont,
    spaceabove=0pt,
    spacebelow=0pt,
    mdframed={style=numexampleframe}
]{numexamplebox}

\declaretheoremstyle[
    headfont=\bfseries\sffamily\color{Maroon},
    bodyfont=\normalfont,
    spaceabove=0pt,
    spacebelow=0pt,
    mdframed={style=numobsframe}
]{numobsbox}

\declaretheoremstyle[
    headfont=\bfseries\sffamily\color{NotationTitle},
    bodyfont=\normalfont,
    spaceabove=0pt,
    spacebelow=0pt,
    mdframed={style=numnotationframe}
]{numnotationbox}

% Environments. Same names as your earlier notes.
\declaretheorem[style=numresultbox, name=Definición, numberwithin=chapter]{definition}
\declaretheorem[style=numdefinitionbox, name=Teorema, numberwithin=chapter]{theorem}
\declaretheorem[style=numdefinitionbox, name=Proposición, numberwithin=chapter]{prop}
\declaretheorem[style=numdefinitionbox, name=Lema, numberwithin=chapter]{lema}
\declaretheorem[style=numdefinitionbox, name=Corolario, numberwithin=chapter]{colorary}
\declaretheorem[style=numexamplebox, numbered=no, name=Ejemplo]{eg}
\declaretheorem[style=numobsbox, numbered=no, name=Observación]{observation}
\declaretheorem[style=numnotationbox, numbered=no, name=Notación]{notation}
\declaretheorem[style=numexamplebox, numbered=no, name=Comentario]{remark}

